% Sources : notes/etat_de_l_art.md, fiches de lecture dans notes/, notes/related_work.md.
% Règle : un chiffre d'un papier n'est écrit que si une fiche existe dans notes/ ; sinon \averifier{}.
\section{État de l'art}\label{sec:etat-art}
	\subsection{Inférence d'appartenance sur les grands modèles de langage}\label{sec:ea-mia}
		\afaire{Loss, zlib, modèle de référence, comparaison avec des voisins, Min-K\%, Min-K\%++.}
	\subsection{Limites méthodologiques des évaluations}\label{sec:ea-critique}
		\afaire{Biais de distribution temporelle, baselines aveugles, passage au niveau jeu de données.}
	\subsection{Travaux spécifiques au code source}\label{sec:ea-code}
		\afaire{Mémorisation, MIA sur modèles de complétion, perturbations sémantiquement équivalentes.}
	\subsection{Accès boîte noire et API fermées}\label{sec:ea-boite-noire}
		\afaire{Méthodes sans log-probabilités du prompt.}
	\subsection{Approches proactives : tatouage du code}\label{sec:ea-tatouage}
		\afaire{CoProtector, CodeMark et leurs limites.}
	\subsection{Positionnement de CodeTrace}\label{sec:ea-positionnement}
		\afaire{Lacune visée et différence avec les travaux les plus proches.}
